\documentstyle[% 


righttag% 


,11pt% 


,amssymb% 


,russian%               remove in Eng. 


,izvru%                 Set izven% in Eng. 


% ,izvru-ks             for brief communications 


]{amsart} 


 


% 


 


\newcommand{\la}{\langle} 


\newcommand{\ra}{\rangle} 


\newcommand{\ul}{\underline} 


\newcommand{\ad}{\operatorname{ad}} 


\newcommand{\Ker}{\operatorname{Ker}} 


\newcommand{\Hom}{\operatorname{Hom}} 


\newcommand{\tr}{\operatorname{tr}} 


\newcommand{\tts}[1]{} 


\renewcommand{\baselinestretch}{0.98}


 


\theoremstyle{plain} 


\theorembodyfont{\it} 


\newtheorem{propnonum}{\hskip\parindent Предложение} 


\renewcommand{\thepropnonum}{} 


 


%\hoffset=-15mm%                  For Eng. 


\setcounter{page}{27} 


\god{2002} 


\nomer{11~(486)} %                        Remove in Eng. 


%\nomer{Vol.\,46, No.\,11}%               Set in Eng. 


%\str{\pageref{begin}--\pageref{end}}%  Set in Eng. 


\udk{514.76} 


 


\author{\it А.А.\,Ермолицкий} 


% NB:   ^^ remove in Eng. 


\title{Теорема о шляпе и проблемы классификации  


структур на римановых многообразиях} 


 


\date{} 


%\pagestyle{myheadings}%         Set in Eng. 


\pagestyle{plain} %              Remove in Eng. 


\begin{document} 


%\thispagestyle{plain}%          Set in Eng. 


\maketitle 


\label{begin} 


 


 


В работе рассматривается теорема о шляпе, позволяющая производить 


изгибание 


подрасслоения в главном $G$-расслоении над дифференцируемым многообразием. 


Далее с использованием ранее введенного автором понятия 


второго фундаментального 


тензорного поля $G$-структуры на римановом многообразии изучаются вопросы 


классификации таких структур. Получены теорема об ``алгебраизации'' 


проблемы 


классификации, а также пример почти эрмитовой структуры, принадлежащей 


разным классам в разных точках многообразия. 


 


\section{Теорема о шляпе и изгибание подрасслоения в главном 


$G$-расслоении} 


 


$\bold 1^\circ$. Пусть $M$ --- связное многообразие, $P(M,G,\pi)$ --- 


главное расслоение 


над $M$. На многообразии $M$ возьмем риманову метрику $g$ и для точки 


$p\in M$ 


рассмотрим геодезические шары $B(p;R/2)\subset B(p;R)\subset U$, где $U$ 


--- 


локальная карта $P(M,G,\pi)$. Таким образом, 


существует диффеоморфизм $\psi:\pi^{-1}(U)\to U\times G:u\mapsto 


(\pi(u),\varphi(u))$, 


где $\varphi:\pi^{-1}(U)\mapsto G$, $\varphi(u a)=\varphi(u)a$ для всех 


$u\in\pi^{-1}(U)$ и $a\in G$. 


 


Далее, если $\psi^{-1}:U\times G\to\pi^{-1}(U)$ --- обратное отображение, 


то $\psi^{-1}|U\times\{e\}$ 


определяет сечение $s_1:U\to\pi^{-1}(U)$. Рассмотрим 


сечение $s_2:U\to\pi^{-1}(U)$ такое, что $s_1(p)=s_2(p)$. Пусть


$v:U\to G$, $x\mapsto v(x)$, --- отображение, определяемое соотношением


$\varphi(s_2(x))=\varphi(s_1(x)v(x))$. 


Если $\exp_p$ --- экспоненциальное отображение pимановой связности 


$\nabla$ в точке $p$, то отображение


$\exp^{-1}_p$ определено на $\ov B(p;R)$ и можно определить следующее 


сечение $s$ над $U$: 


\begin{equation} 


s(x)= 


\begin{cases} 


s_1(x), & x\in U\setminus\ov B(p;R); \\ 


s_2(x)\left(v\left[\exp_p\left(\frac{2t-R}{t} 


\exp^{-1}_p(x)\right)\right]\right)^{-1}, & x\in\ov B(p;R) 


\setminus\ov B(p;R/2); \\ 


s_2(x), & x\in\ov B(p;R/2), 


\end{cases} 


\end{equation} 


где $\exp^{-1}_p(x)=t\xi$, $\|\xi\|=1$. В результате получена


 


\begin{thm}[{\series{m}\shape{n}\selectfont о шляпе}] Если $s_1$, $s_2$ 


--- локальные сечения главного 


расслоения $P(M,G,\pi)$ и $s_1(p)=s_2(p)$, то 


существуют последовательность 


окрестностей $U\supset B(p;R)\supset B(p;R/2)\ni p$ 


и локальное сечение $s:U\to P(M,G,\pi)$ такие, что $s=s_1$ 


на $U\setminus\ov B(p;R)$ и $s=s_2$ на $\ov B(p;R/2)$. 


Такое сечение может задаваться, 


например, формулой 


$(1)$. 


\end{thm} 


 


Далее, пусть $P'(M,G',\pi)$ --- редуцированное подрасслоение 


расслоения $P(M,G,\pi)$, 


$\{U_\alpha\}$ --- открытое покрытие на $M$ с множеством функций перехода 


$\psi_{\beta\alpha}$, 


принимающих значения из $G'$, а $U\ni p$ --- одна из окрестностей этого 


покрытия. 


 


Рассмотрим многообразие $M'=M\setminus\ov B(p;R)$ с соответствующим 


открытым 


покрытием $\{U'_\alpha\}$, $U'_\alpha=U_\alpha\setminus\ov B(p;R)$ и 


функциями перехода $\psi'_{\beta\alpha}=\psi_{\beta\alpha|M'}$, 


тогда открытое покрытие $\{U'_\alpha;U\}$ с функциями перехода 


$\psi'_{\beta\alpha}$ и 


сечением $s$ над $U$, 


определенным формулой (1), задает некоторое новое подрасслоение 


расслоения $P(M,G,\pi)$. Таким образом, получена


 


\begin{thm} 


Пусть $P'(M,G',\pi)$ --- редуцированное подрасслоение 


расслоения $P(M,G,\pi)$, $\{U_\alpha\}$ --- открытое покрытие 


на $M$, $s_\alpha:U_\alpha\to P'(M,G',\pi)$ --- соответствующие 


сечения, $p\in U_1$ и $s_1(p)=\ov s_2(p)$, где $\ov s_2:U_1\to P(M,G,\pi)$ 


--- некоторое сечение. Тогда существует такое подрасслоение $\ov 


P{}'(M,G',\pi)$ 


расслоения $P(M,G,\pi)$, что сечение


$s:U_1\to\pi^{-1}(U_1)$, определяемое формулой $(1)$,


является сечением подрасслоения $\ov P{}'(M,G',\pi)$. 


\end{thm} 


 


\begin{defn} 


Подрасслоение $\ov P{}'(M,G',\pi)$ будем называть 


изгибанием в $P(M,G,\pi)$ подрасслоения $P'(M,G',\pi)$ по формуле (1) в 


точке $p$. 


\end{defn} 


 


$\bold 2^\circ$. Напомним, что связность $\Gamma$ в расслоении 


$P(M,G,\pi)$ --- это 


сопоставление подпространства $Q_u$ из $T_u(P)$ каждой точке $u\in P$ 


такое, что 


\begin{itemize} 


\item[(a)] $T_u(P)=G_u\oplus Q_u$, где $G_u$ --- подпространство из 


$T_u(P)$, касательное к слою, проходящему через $u$;


\item[(b)] $Q_{u a}=(R_a)_*Q_u$, $u\in P$, $a\in G$, $R_a u=u a$;


\item[(c)] $Q_u$ зависит дифференцируемо от $u$. 


\end{itemize} 


 


\begin{thm} 


Пусть для некоторого $u_0\in P(M,G,\pi)$ задано некоторое


подпространство $Q_{u_0}$ из $T_{u_0}(P)$ такое, что 


$T_{u_0}(P)=G_{u_0}\oplus Q_{u_0}$ и $P'(M,G',\pi)$ --- 


редуцированное подрасслоение расслоения $P(M,G,\pi)$. Тогда существуют 


подрасслоение $\ov P{}'(M,G',\pi)$ и связность $\Gamma$ в расслоении 


$P(M,G,\pi)$ такие, что 


 


$1)$ связность $\Gamma$ в $P$ редуцируема к связности $\Gamma'$ в 


$\ov P{}'$, 


 


$2)$ $Q_{u_0}$ есть горизонтальное подпространство в $T_{u_0}(P)$ 


связности $\Gamma$. 


\end{thm} 


 


\begin{pf} 


Пусть $p=\pi(u_0)$, тогда для любого 


$u=u_0a\in\pi^{-1}(p)$ определим $Q_u=(R_a)_*Q_{u_0}$. Если $U\ni p$ --- 


локальная 


карта подрасслоения $P'(M,G',\pi)$ с 


соответствующим сечением $s_1:U\to \pi^{-1}(U)$, $s_1(p)=u_1$, 


то $\pi^{-1}(U)$ диффеоморфно $U\times G$, 


и, выбрав римановы метрики на $U$ и $G$ соответственно, можно 


ввести риманову 


метрику прямого произведения $\wt g$ на $\pi^{-1}(U)$, считая 


диффеоморфизм $\psi^{-1}$ 


изометрией. Взяв геодезический шар $B(u_1)=\wt\exp_{u_1}(V)$, 


где $V\subset T_{u_1}(P)$, можно определить 


подмногообразие $\wt\exp_{u_1}(V\cap Q_{u_1})$ в $B(u_1)$ такое, что $\pi$ 


есть диффеоморфизм 


на $\wt\exp_{u_1}(V\cap Q_{u_1})$ и касательное подпространство к этому 


подмногообразию 


в точке $u_1$ совпадает с $Q_{u_1}$. Таким образом, определяется 


сечение $s_2:B(p;R)\to P(M,G,\pi)$, $x\mapsto\pi^{-1}(x)\cap 


\wt\exp_{u_1}(V\cap Q_{u_1})$, 


где $B(p;R)=\pi(B(u_1))$. 


 


Используя теорему 2, построим по формуле (1) изгибание $\ov P{}'(M,G',\pi)$


подрасслоения\linebreak $P'(M,G',\pi)$. Далее, для каждого 


$u=s_2(x)$, $x\in B(p;R/2)$, определим горизонтальное подпространство 


$Q_u$ как касательное пространство к подмногообразию 


$s_2(B(p;R/2))$, а в остальных точках из $\pi^{-1}(B(p;R/2))$


определим горизонтальное подпространство, исходя из правой 


инвариантности. Таким образом, получили главное расслоение


$\ov P{}'(M,G',\pi)$ и связность $\Gamma'$ в расслоении


$P{}'(M,G',\pi)$, определенную над $B(p;R/2)$. 


Тогда эта связность может быть продолжена до связности в $\ov 


P{}'(M,G',\pi)$ 


над $M$, если $M$ паракомпактно [1], и из правой инвариантности, до 


связности $\Gamma$ в $P(M,G,\pi)$ над $M$. Горизонтальное подпространство 


в $T_{u_0}(P)$ связности $\Gamma$ 


совпадает с $Q_{u_0}$ по построению. 


\end{pf} 


 


 


\section{Проблемы классификации $G$-структур на римановых многообразиях} 


 


$\bold 1^\circ$. Пусть $L(M)$ --- расслоение линейных реперов над 


многообразием 


$M$, $P(G)$ --- $G$-структура над $M$, допускающая редукцию к максимальной 


компактной подгруппе $H$ группы $G$, $H=G\cap O(n)$. Расширяя $P(H)$ до 


группы $O(n)$, получим подрасслоение ортонормальных реперов $O(M)$, 


определяющее специальную риманову метрику $g=\la,\ra$ структуры $P(G)$ 


[2]. 


Для алгебр Ли $\ul o$, $\ul h$ групп Ли $O(n)$ и $H$


имеем $\ul o=\ul h\oplus \ul m$, где $\ul m$ --- ортогональное 


дополнение относительно формы Киллинга, $\ad(H)\ul m=\ul m$ из 


биинвариантности 


формы Киллинга. Если $\omega$ --- $o$-значная форма римановой связности 


в $O(M)$, то $\ov\omega=\omega_{|\ul h}$ определяет некоторую связность в 


$P(H)$. 


Связности $\omega$, $\ov\omega$ могут быть 


расширены до линейных связностей с соответствующими ковариантными 


производными $\nabla$, $\ov\nabla$. 


 


\begin{defn}[{\series{m}\shape{n}\selectfont [3]}] 


Тензорное поле $h=\nabla-\ov\nabla$ называется вторым 


фундаментальным тензорным полем структуры $(P(G),g)$; связность 


$\ov\nabla$ 


называется канонической связностью; пара $(P(G),g)$ называется 


квазиоднородной структурой, если $\ov\nabla h=0$ на $M$. 


\end{defn} 


 


Отметим, что понятие квазиоднородности рассматривалось в [4]. 


 


Пусть $\wt g$ --- произвольная риманова метрика на


множестве $M$, $X,Y\in\frak X(M)$. 


Следующие примеры можно найти в [2]:


 


1) $(P,g)$, $P^2=I$, $g(X,Y)=\wt g(X,Y)+\wt g(P X,P Y)$, $\ov\nabla_X 


Y=\frac{1}{2}(\nabla_X Y+P\nabla_X P Y)$,


 


2) $(J,g)$, $J^2=-I$, $g(X,Y)=\wt g(X,Y)+\wt g(J X,J Y)$, $\ov\nabla_X 


Y=\frac{1}{2}(\nabla_X Y-J\nabla_X J Y)$,


 


3) $(F,g)$, $F^3+F=0$, $\ov\nabla_X Y=\nabla_X Y-\frac{1}{2}F \nabla_X F 


Y+\nabla_X F^2Y+F^2\nabla_X Y+\frac{3}{2}F^2\nabla_X F^2Y$,


 


4) $(S,g)$, $S^k=I$, $g(X,Y)=\sum\limits^{k-1}_{j=0}\wt g(S^j X, S^j Y)$, 


$\ov\nabla_X Y=\frac{1}{k}\sum\limits^{k-1}_{j=0}S^j \nabla_X S^{k-j}Y$. 


 


Следующие примеры приводятся впервые: 


 


5) $(P_1,P_2,g)$, $P^2_1=P^2_2=I$, $P_1P_2=-P_2P_1=J$, $g(X,Y)=\wt 


g(X,Y)+\wt g(P_1X,P_1Y)+\wt g(P_2X,P_2Y)+\wt g(J X,J Y)$, $\ov\nabla_X 


Y=\frac{1}{4}(\nabla_X Y+P_1\nabla_X P_1Y+P_2 \nabla_X P_2Y-J \nabla_X J 


Y)$,


 


6) $(J_1,J_2,g)$, $J^2_1=J^2_2=-I$, $J_1J_2=-J_2J_1=J_3$, $g(X,Y)=\wt 


g(X,Y)+\wt g(J_1X,J_1Y)+\wt g(J_2X,J_2Y)+\wt g(J_3X,J_3Y)$, $\ov\nabla_X 


Y=\frac{1}{4}(\nabla_X Y-J_1\nabla_X J_1Y-J_2\nabla_X J_2Y-J_3\nabla_X 


J_3Y)$. 


\medskip 


 


$\bold 2^\circ$. Напомним, что для каждого $\xi\in\bold R^n$ и $u\in L(M)$ 


существует единственный вектор $(B(\xi))_u$ в $Q_u$ такой, что 


$\pi((B(\xi))_u)=u(\xi)$. \ $B(\xi)$ называется стандартным горизонтальным 


векторным полем, соответствующим $\xi\in\bold R^n$. Если $A^*$ --- 


фундаментальное 


векторное поле, соответствующее $A\in\ul o$, то на $O(M)$ определяется 


риманова метрика $g_\nabla$: 


 


$g_\nabla(B(\xi),B(\eta))=\la\xi,\eta\ra$, $\xi,\eta\in\bold R^n$; 


 


$g_\nabla(A^*_1,A^*_2)=-\tr(A_1A_2)$, $A_1,A_2\in\ul o$; 


 


$g_\nabla(B(\xi),A^*)=0$, $\xi\in\bold R^n$, $A\in\ul o$. 


 


Пусть $Q=\Ker\omega$, $\ov Q=\Ker\ov\omega$. Рассмотрим 


распределения $V^{\ul h}=\{A^*:A\in\ul h\}$, $V^{\ul m}=\{A^*:A\in\ul 


m\}$. 


Тогда очевидно, что $V^{\ul m}\bot V^{\ul h}$ и $T_u P(H)\cap V_u=V^{\ul 


h}_u$, 


где $u\in P(H)$, $V_u$ --- касательное пространство к слою в $O(M)$. 


 


\begin{propnonum} 


Горизонтальное распределение $\ov Q_u=\Ker\ov\omega$ 


канонической связности структуры $(P(H),g)$ определяется по формуле 


$$ 


u\to\ov Q_u=T_u(P(H))\cap V^{\ul h^\bot}_u= 


T_u(P(H))\cap(V^{\ul m}_u\oplus Q_u),\quad u\in P(H). 


$$ 


\end{propnonum} 


 


\begin{pf} 


Так как $\omega(\ov Q_u)\subset\ul m$, то $\ov\omega(\ov Q_u)=0$, 


$T_u(P(H))=V^{\ul h}_u\oplus\ov Q_u$, 


следовательно, $\dim\ov Q_u=n$ и $\ov Q_u=\Ker\ov\omega$. 


\end{pf} 


 


Для некоторого $u_0\in O(M)$, $\pi(u_0)=p$, рассмотрим произвольное 


подпространство $\ov Q_{u_0}$ размерности $n$ из $V^{\ul h}_u\oplus 


Q_{u_0}$ такое, что $\ov Q_{u_0}\cap V^{\ul m}_{u_0}=0$. 


Тогда $T_{u_0}(O(M))=V_{u_0}\oplus\ov Q_{u_0}$ 


и по теореме~3 существуют такая структура $P'(H)$ и 


связность $\ov\Gamma{}'$ в $P'(H)$, что $\ov Q_{u_0}$ есть горизонтальное 


подпространство в $T_{u_0}(P)$ связности $\ov\Gamma{}'$. Пусть


$\ov\Gamma$ --- каноническая связность структуры $P'(H)$ из определения 1. 


По доказанному выше предложению $\ov Q_{u_0}$ является


горизонтальным подпространством связности $\ov\Gamma$ в $T_{u_0}(P)$. 


 


Далее, пусть $\varphi\in\Hom(\bold R^n,\ul m)$ и 


$\xi_1,\dotsc,\xi_n\in\bold R^n$ --- базисные 


векторы, $V_i=(\varphi(\xi_i))^*_{u_0}+(B(\xi_i))_{u_0}$. Тогда 


определяется линейная 


оболочка $\ov Q_{u_0}=[V_1,\dotsc,V_n]$ такая, что 


$T_{u_0}(O(M))=V_{u_0}\oplus\ov Q_{u_0}$. Обратно, 


для такого подпространства $\ov Q_{u_0}$ линейное отображение $\varphi$ 


определяется 


формулой $\varphi=\omega_{u_0}\circ\ov B$, где $(\ov B(\xi))_{u_0}$ --- 


единственный вектор в $\ov Q_{u_0}$ 


такой, что $\pi(\ov B(\xi)_{u_0})=u_0(\xi)$. 


 


Таким образом, из вышесказанного следует 


 


\begin{thm} 


Пусть на многообразии $M$ существует 


структура $P(H)$, $H\subset O(n)$, $p\in M$, и $\{P(H)\}$ --- 


множество всевозможных $H$-структур в 


расслоении $O(M)$, которое получено расширением $H$ до $O(n)$. 


Тогда множество горизонтальных распределений в точке $p$ канонических 


связностей этих $H$-структур находится во взаимно однозначном 


соответствии с $\Hom(\bold R^n,\ul m)$. 


\end{thm} 


 


Для точки $p=\pi(u_0)$ рассмотрим $T=T_p(M)$, $X,Y,Z\in T$, 


$h=\nabla-\ov\nabla$, $\tau=\omega-\ov\omega$, где 


$\ov\omega$ --- форма связности, а $\ov\nabla$ ---


ковариантная производная канонической 


связности $H$-структуры $P(H)$, 


имеющей в $u_0$ горизонтальное распределение $\ov Q_{u_0}$. 


Тогда $h_X Y=u_0\tau(B(\xi))u^{-1}_0Y$, где $u_0\xi=X$ [5]. Так как 


$u_0:\bold R^n\to T$ есть изометрия, то 


$$ 


h_{X Y Z}=\la h_X Y,Z\ra=\la u_0\tau(B(\xi))u^{-1}_0Y,Z\ra= 


\la-\ov\omega(B(\xi))u^{-1}_0Y,u^{-1}_0Z\ra= 


-\la\ov\omega(B(\xi))\eta,v\ra= 


\la\omega(\ov B(\xi))\eta,v\ra, 


$$ 


где $u_0\eta=Y$, $u_0v=Z$. 


 


Для множества $\{P(H)\}$ всевозможных $H$-структур в расслоении $O(M)$ 


рассмотрим 


подмножество $\ov{\bold T}(T)$ в $\overset{3}{\oplus}{}T^*$,


состоящее из всех тензоров 


$h_p$, где $h$ --- второе фундаментальное 


тензорное 


поле $P(H)$. Определяется взаимно однозначное 


соответствие 


$$ 


\Phi:\ov{\bold T}(T)\to\Hom(\bold R^n,\ul m),


\ \ h_p\mapsto \varphi=\omega_{u_0}\circ\ov B= 


-\ov\omega_{u_0}\circ B 


$$ 


($h_p$ и $\ov\omega_{u_0}$ определяют друг друга взаимно однозначно). 


 


Нетрудно убедиться, что отображение $\Phi$ линейно, и мы имеем линейный 


изоморфизм между $\ov{\bold T}(T)$ и $\Hom(\bold R^n,\ul m)$. 


\medskip 


 


$\bold 3^\circ$. Риманова метрика $g$ превращает пространство $T=T_p(M)$ 


в евклидово 


пространство. Если $E_1,\dotsc,E_n$ --- произвольный ортонормальный 


базис $T$, то $\ov{\bold T}(T)$ является евклидовым пространством 


относительно следующего скалярного произведения: 


$$ 


\la h,h'\ra=\sum_{i,j,k}h_{E_i E_j E_k}\cdot 


h'_{E_i E_j E_k},\quad h,h'\in\ov{\bold T}(T). 


$$ 


Естественное действие группы $H$ на $T$ индуцирует действие


этой группы на $\ov{\bold T}(T)$ по формуле 


$$ 


(a h)_{X Y Z}=h_{a^{-1}X a^{-1}Y a^{-1}Z},\quad 


X,Y,Z\in T,\ \ a\in H. 


$$ 


 


Пусть $\ov{\bold T}(T)=\underset{\alpha\in A}{\oplus}{}\ov{\bold 


T}_\alpha$ --- разложение пространства $\ov{\bold T}(T)$ на инвариантные и 


неприводимые компоненты относительно действия группы $H$. Заметим, что 


изоморфизм $\Phi$ индуцирует скалярное произведение на $\Hom(\bold R^n,\ul 


m)$. 


 


\begin{thm} 


Изоморфизм $\Phi$ осуществляет взаимно однозначное соответствие 


между\linebreak $\{\ov{\bold T}_\alpha\}_{\alpha\in A}$ и множеством всех 


инвариантных и 


неприводимых компонент пространства\linebreak


$\Hom(\bold R^n,\ul m)$ относительно следующего действия 


группы $H{:}$ 


\begin{equation} 


H\times\Hom(\bold R^n,\ul m)\to\Hom(\bold R^n,\ul m),


\ \ \varphi\mapsto a\varphi,\ \ 


(a\varphi)\xi=\ad(a)\varphi(a^{-1}\xi),\ \ a\in H, 


\ \ \xi\in\bold R^n. 


\end{equation} 


\end{thm} 


 


\begin{pf} 


Так как $h_{X Y Z}$ не зависит от выбора репера $u\in O(M)$, то 


$\la\ov\omega_u(B(\xi))\eta,v\ra= \la\ov\omega_{u 


a}(B(a^{-1}\xi))a^{-1}\eta,a^{-1}v\ra$, 


где $u\xi=X$, $u\eta=Y$, $u v=Z$. 


\begin{multline*} 


(a h)_{X Y Z}=h_{a^{-1}X a^{-1}Y a^{-1}Z}= 


-\la\ov\omega_u(B(a^{-1}\xi))u^{-1}a^{-1}Y,u^{-1}a^{-1}Z\ra= 


-\la\ov\omega_u(B(a^{-1}\xi))a^{-1}\eta,a^{-1}v\ra= \\ 


% 


=-\la\ov\omega_{u a^{-1}}(B(\xi))\eta,v\ra= 


-\la\ov\omega_{u a^{-1}}(R_{a^{-1}_*}R_{a_*}B(\xi)) 


\eta,v\ra=-\la\ad(a)\ov\omega_u(B(a^{-1}\xi))\eta,v\ra. 


\end{multline*} 


Если $\Phi(h_p)=\varphi$, то $\Phi(a h_p)=a\varphi=a\Phi(h_p)$, где 


$a\varphi$ определяется из (2). 


Действие группы $H$ 


на $\Hom(\bold R^n,\ul m)$ действительно определяется 


из (2), т.\,к. $\ad(H)\ul m=\ul m$ и 


$((a b)\varphi)\xi=\ad(a b)\varphi(b^{-1}a^{-1}\xi)= 


\ad(a)\ad(b)\varphi(b^{-1}a^{-1}\xi)= \ad(a)((b\varphi)a^{-1}\xi)= 


(a(b\varphi))\xi$. Так как $\Phi(a h_p)=a\Phi(h_p)$, то $\Phi$ является 


изоморфизмом $\Hom(\bold R^n,\ul m)$ и $\ov{\bold T}(T)$ как 


$H$-пространств и дает взаимно однозначное соответствие между $\{\ov{\bold 


T}_\alpha\}_{\alpha\in A}$ 


и неприводимыми компонентами $\Hom(\bold R^n,\ul m)$. 


\end{pf}


 


$\bold 4^\circ$. Отметим, что введеннoe автором тензорноe полe $h$


применялось в [2] к 


проблемам классификации структур на римановых многообразиях. 


В частности, в терминах поля $h$ в [2] была получена 


классификация почти эрмитовых структур ($H=U(n)$, всего 16 классов), 


изоморфная классификации в [6], а также классификация почти контактных 


метрических структур ($H=U(n)\times1$, всего $2^{12}$ классов), изоморфная 


классификации в [7], [8]. 


 


Аналогичным образом,


с помощью поля $h$ может быть переписана и классификация 


из [9]. Все перечисленные классификации фактически проводятся 


инфинитезимально (в точке $p\in M$). Естественно дать следующее 


 


\begin{defn} 


Будем говорить, что структура $(P(H),g)$ принадлежит 


классу $\ov{\bold T}_\alpha$, $\alpha\in A$, на $M$, если $h_p\in\ov{\bold 


T}_\alpha(T_p(M))$ для всех $p\in M$. 


\end{defn} 


 


Если тензор $h_p\in\ov{\bold T}_\alpha(T)$, $p=\pi(u_0)$, то единственным 


образом получим 


гомоморфизм $\varphi=\Phi(h_p)$, который определяет 


подпространство $\ov Q_{u_0}\subset V^{\ul m}_{u_0}\oplus Q_{u_0}$ такое, 


что $T_{u_0}(O(M))=V_{u_0}\oplus\ov Q_{u_0}$. Из доказательства теоремы 4


видно, 


что существует структура $P'(H)\subset O(M)$ такая, что $\ov Q_{u_0}$ --- 


горизонтальное подпространство ее канонической связности $\ov\Gamma$, 


т.\,е. $h_p$ является вторым фундаментальным тензором структуры 


$(P'(H),g)$ 


в точке $p$. Беря в различных точках $p_1,\dotsc,p_s$ тензоры 


$h_{p_1},\dotsc,h_{p_s}$, 


принадлежащие различным классам $\ov{\bold T}_\alpha$, $\alpha\in A$, и, 


проводя изгибание 


структуры $P(H)$ в $O(M)$ в каждой из указанных точек, получим 


$H$-структуру $\ov P(H)$ на $M$, для которой второе фундаментальное 


тензорное 


поле $h$ удовлетворяет условиям $h(p_i)=h_{p_i}$, $i=1,\dotsc,s$. 


В частности, для почти эрмитовой структуры $(J,g)$, $J^2=-I$, $H=U(n)$ 


на многообразии $M$ можно выбрать 16 точек и почти эрмитову структуру 


$(\ov J,g)$ так, что $h_{p_i}\in\ov{\bold T}_i$, $i=0,1,\dotsc,15$, где $h$ 


--- второе 


фундаментальное поле пары $(\ov J,g)$ (см. пример 2). Это значит, что 


структура $(\ov J,g)$ будет кэлеровой в $p_0$, приближенно


(nearly)-кэлеровой в $p_1$, 


почти (almost)-кэлеровой в $p_2$, интегрируемой в $p_6$ и т.д. 


 


Таким образом, принадлежность структуры к определенному классу в 


общем случае не сохраняется. 


 


\begin{thm}[{\series{m}\shape{n}\selectfont [2]}] 


Пусть почти контактная метрическая 


структура $($почти эрмитова структура$)$ на множестве $M$ является 


квазиоднородной структурой $(\ov\nabla h=0)$, принадлежащей 


классу $\ov{\bold T}_\alpha$ для некоторой точки $p\in M$. Тогда эта 


структура будет 


из класса $\ov{\bold T}_\alpha$ для каждой точки многообразия $M$. 


\end{thm} 


 


{\bf Проблема.} {\it Остается ли эта теорема справедливой в общем 


случае для структуры $(P(H),g)$ на множестве $M$, $H\subset O(n)$}? 


 


 


\begin{thebibliography}{9} 


 


\bibitem{1} 


Кобаяси Ш., Номидзу К. {\em Основы дифференциальной 


геометрии}. Т.\,1. -- М.: Наука, 1981. -- 344~с. 


\bibitem{2} 


Ермолицкий А.А. {\em Римановы многообразия с 


геометрическими структурами}. -- Минск: БГПУ, 1998. -- 195 с. 


\bibitem{3} 


Ермолицкий А.А. {\em Вторая квадратичная форма $G$-структуры 


и почти особые структуры} // 


Докл. Болг. АН. -- 1981. -- Т.\,7. -- \No\,7. -- С.\,963--964. 


\bibitem{4} 


Кириченко В.Ф. {\em Квазиоднородные многообразия и обобщенные почти 


эрмитовы структуры} // 


Изв. АН СССР. -- 1983. -- Т.\,47. -- \No\,6. -- С.\,1208--1223. 


\bibitem{5} 


Бишоп Р.Л., Криттенден Р.Д. 


{\em Геометрия многообразий}. -- М.: Мир, 1967. -- 335 с. 


\bibitem{6} 


Gray A., Hervella L.M. {\em The sixteen classes of almost 


Hermitian manifolds and their linear 


invariants} // Ann. mat. pura appl. -- 1980. -- V.\,123. -- P.\,35--58. 


\bibitem{7} 


Alexiev V.A., Ganchev G.T. {\em On the classification of the almost 


contact metric manifolds} 


// Proceedings of the 15$^{\mathrm{th}}$ spring conference of the Union of 


Bulgarian Mathematicians. Math. 


and Educ. in Math. -- 1986.-- P.\,155--161. 


\bibitem{8} 


Chinea D., Gonzalez C. {\em A classification of almost contact 


metric manifolds} // Ann. mat. pura appl. 


-- 1990. -- V.\,156. -- P.\,15--36. 


\bibitem{9} 


Naveira A.M. {\em A classification of Riemannian almost-product 


manifolds} // Rend. math. appl. -- 1983. -- V.\,3. 


-- \No\,3. -- P.\,577--592. 


 


\end{thebibliography} 


 


\vskip 1cm 


 


\post{\it Белорусский институт правоведения}{\it Поступила} 


\post{}{25.10.2001} 


 


\label{end} 


\end{document} 


